\chapter{\textit{preferred}-Unabhängigkeit}
\label{ch:preferred}

Einleitung in Kapitel 

\section{Komplexität}

\section{Aussagenlogische Kodierung der \textit{preferred}-Semantik}
Wir erinnern zuerst die Definition der \textit{preferred}-Semantik. Um die nachfolgende Kodierung 
zu vereinfachen, 
nutzen wir als Ausgangspunkt die \textit{complete}-Semantik. Sei \(F=(A,R)\) ein AF 
und seien \(S, S' \subseteq A\) \textit{complete}-Extensionen von \(F\). $S$ ist eine \textit{preferred}-Extension
genau dann, wenn für jedes $S' \subseteq S$ gilt.

Um die \textit{preferred}-Semantik zu definieren müssen wir  eine Aussage über die Menge aller \textit{complete}-
Extensionen treffen. 

\begin{Definition}
    \label{def:geq-pl}
    Sei \(F=(A,R)\) ein AF  und \(I \text{ und } I'\) zwei Kopien eines aussagenlogischen Labellings von $F$. Wir definieren
    \[
    I \geq I'
    := 
    \bigwedge_{a \in A}
    (
        i'_a \rightarrow i_a
    )
    \]
\end{Definition}

\begin{Lemma}
    \label{lem:geq-pl-correct}
    Sei 

    Dann gilt: 

\end{Lemma}

Wir nutzen die in \ref{def:vars} eingeführten aussagenlogischen Variablen und die in \ref{def:pl-one} definierte
Formel, um die \textit{preferred}-Labellings aussagenlogisch zu kodieren. 

\begin{Definition}
    \label{def:prf-pl}
    Sei \(F = (A,R)\) ein AF und seien \(I, O, U\) wie zuvor definiert. 
    Eine aussagenlogische Kodierung eines \textit{preferred}-Labellings lautet
    \[
    \mathrm{prf}_F(\mathcal{L})
    :=
    \forall \mathcal{L'}
    \bigl(
        \mathrm{co}_F(\mathcal{L}) 
        \land 
        (
            \mathrm{co}_F(\mathcal{L'})
            \rightarrow
            I \geq I'
        ) 
    \bigr)
    \]
\end{Definition}

\begin{Lemma}
    \label{lem:prf-pl-correct}
\end{Lemma}

\begin{proof}
    \label{pr:prf-pl-correct}
\end{proof}

\section{QBF-Kodierung der \textit{preferred}-Unabhängigkeit}
